\section{Results}
\label{sec:results}

We report results across five Gowalla state splits. The story splits in two: the substrate's effect is asymmetric across the two tasks (\S\ref{subsec:results-substrate}), and joint multi-task training over Check2HGI adds a small category lift while paying the textbook region cost (\S\ref{subsec:results-mtl-stl}). Section~\ref{subsec:results-deltam} reports the joint $\Delta_m$ score; Section~\ref{subsec:results-external} situates the model against external baselines.

\subsection{Substrate ablation: the embedding is asymmetric across tasks}
\label{subsec:results-substrate}

On next-category, per-visit context lifts macro-F1 by $+14.18$ to $+29.02$ percentage points (pp) at every state (smallest at AZ, largest at FL). Table~\ref{tab:substrate}(a) reports the numbers under the matched-head GRU single-task ceiling, with substrate as the only changing factor. The paired Wilcoxon signed-rank test reaches the $n=5$ ceiling everywhere ($p=0.0312$, 5/5 folds positive). Check2HGI macro-F1 is stable across four extra seeds $\{0, 1, 7, 100\}$ at AL/AZ/FL with seed-$\sigma\leq 0.17$ pp. Two robustness probes back this: the AL/AZ gap is head-invariant across linear, single-layer, GRU, and LSTM probes (in the anonymous code release), and the FL/CA/TX gap replicates under the matched-head ceiling. We read the lift as evidence that per-visit context, not POI-level attribute encoding, is what the next-category head exploits.

On next-region, by contrast, per-place HGI ties or marginally exceeds Check2HGI. Table~\ref{tab:substrate}(b) reports Acc@10 under the matched-head STAN-based single-task ceiling. HGI is nominally above Check2HGI by 1.6 to 3.1 pp at every state; on FL/CA/TX the gap shrinks to 1.6 to 2.1 pp and TOST non-inferiority at $\delta=2$ pp passes at CA and TX. At FL the $\delta=3$ pp test passes. AL and AZ stay outside both margins. The short-horizon transition structure driving the region head is captured equally well, or fractionally better, by per-place embeddings whose category-level pooling smooths the per-visit variance that helps the category head.

\begin{table}[t]
\centering
\caption{Substrate ablation under matched-head single-task ceilings (5-fold CV, seed-42). Cells are mean $\pm$ fold-$\sigma$; $\Delta=$ Check2HGI $-$ HGI; paired Wilcoxon reaches $p=0.0312$ when 5/5 folds agree. (a) Next-category macro-F1 with a GRU head; asterisk marks 5/5 folds positive. (b) Next-region Acc@10 with a STAN-based head + additive log-transition prior; TOST verdicts at $\delta\in\{2,3\}$ pp test non-inferiority (\checkmark) or not ($\times$).}
\label{tab:substrate}
\footnotesize
\renewcommand{\arraystretch}{1.1}

\textbf{(a) Next-category macro-F1.}
\begin{tabular}{l r r r r}
\toprule
State & HGI cat F1 & Check2HGI cat F1 & $\Delta$ (pp) & Wilcoxon $p$ \\
\midrule
Alabama (AL)    & 25.26 $\pm$ 1.18 & 40.76 $\pm$ 1.50 & $+15.50$\textsuperscript{*} & 0.0312 \\
Arizona (AZ)    & 28.99 $\pm$ 0.51 & 43.17 $\pm$ 0.28 & $+14.18$\textsuperscript{*} & 0.0312 \\
Florida (FL)    & 34.41 $\pm$ 1.05 & 63.43 $\pm$ 0.98 & $+29.02$\textsuperscript{*} & 0.0312 \\
California (CA) & 31.13 $\pm$ 1.04 & 59.94 $\pm$ 0.59 & $+28.81$\textsuperscript{*} & 0.0312 \\
Texas (TX)      & 31.89 $\pm$ 0.55 & 60.24 $\pm$ 1.84 & $+28.34$\textsuperscript{*} & 0.0312 \\
\bottomrule
\end{tabular}

\vspace{0.6em}

\textbf{(b) Next-region Acc@10.}
\begin{tabular}{l r r r c c}
\toprule
State & HGI Acc@10 & Check2HGI Acc@10 & $\Delta$ (pp) & TOST $\delta = 2$ & TOST $\delta = 3$ \\
\midrule
Alabama (AL)    & 61.86 $\pm$ 3.29 & 59.15 $\pm$ 3.48 & $-2.71$ & $\times$ & $\times$ \\
Arizona (AZ)    & 53.37 $\pm$ 2.55 & 50.24 $\pm$ 2.51 & $-3.13$ & $\times$ & $\times$ \\
Florida (FL)    & 71.34 $\pm$ 0.64 & 69.22 $\pm$ 0.52 & $-2.12$ & $\times$ & \checkmark \\
California (CA) & 57.77 $\pm$ 1.12 & 55.92 $\pm$ 1.20 & $-1.85$ & \checkmark & \checkmark \\
Texas (TX)      & 60.47 $\pm$ 1.26 & 58.89 $\pm$ 1.28 & $-1.59$ & \checkmark & \checkmark \\
\bottomrule
\end{tabular}
\end{table}


The substrate's effect is asymmetric across the two tasks. The category head reads per-visit context and lifts by an order of magnitude more than its standard-deviation envelope; the region head reads transition structure that is preserved under per-place pooling and ties within $\pm 2$ pp. Section~\ref{sec:mechanism} attributes the asymmetry through a per-visit-context counterfactual replicated at AL and AZ. Any multi-task system on this task pair therefore operates on a substrate that already encodes the easier task more clearly than the harder one.

\subsection{Multi-task vs single-task on the per-visit substrate}
\label{subsec:results-mtl-stl}

Joint multi-task training is positive on next-category at four of five states. Table~\ref{tab:mtl-stl} reports the main contrast: the cross-attention multi-task backbone against the matched-head GRU and STAN-based single-task ceilings on Check2HGI, multi-seed pooled at $n=20$ fold-pairs per state. Category is positive at FL $+1.40$ pp ($p=2\mathrm{e}{-6}$), CA $+1.68$ pp ($p=2\mathrm{e}{-6}$), TX $+1.89$ pp ($p=2\mathrm{e}{-6}$) and AZ $+1.20$ pp ($p<1\mathrm{e}{-4}$). At AL the contrast is statistically significant but small in magnitude ($\Delta=-0.78$ pp, $p=0.036$, $n=20$, under $2\,\%$ relative on the $41\,\%$ F1 scale). The four positive states sit outside the seed and fold envelope; we report the AL negative as is rather than absorbing it.

\begin{table}[t]
\centering
\caption{Headline MTL-vs-STL contrast on the Check2HGI substrate, multi-seed pooled at $n = 20$ fold-pairs per state (seeds $\{0, 1, 7, 100\}$, five user-disjoint folds). The MTL row is the cross-attention multi-task backbone; the STL rows are matched-head single-task ceilings, \texttt{next\_gru} for category and \texttt{next\_stan\_flow} for region. Cells are mean $\pm$ pooled $\sigma$ over $n = 20$ fold-pairs. The $\Delta$ columns report MTL $-$ STL in percentage points; the Wilcoxon $p$-values are paired one-sided at $n = 20$ (alternative ``greater'' for category gains, ``less'' for region drops where the headline direction is negative). FL/CA/TX sit above the rule as the headline scale; AL/AZ sit below as small-scale anchors.}
\label{tab:mtl-stl}
\footnotesize
\resizebox{\textwidth}{!}{
\begin{tabular}{l r r r r r r}
\toprule
& \multicolumn{3}{c}{Next-category (macro-F1)} & \multicolumn{3}{c}{Next-region (Acc@10)} \\
\cmidrule(lr){2-4} \cmidrule(lr){5-7}
State & STL & MTL & $\Delta$ (pp), $p$ & STL & MTL & $\Delta$ (pp), $p$ \\
\midrule
Florida (FL)    & 67.16 $\pm$ 0.13 & 68.56 $\pm$ 0.79 & $+1.40$, $2\mathrm{e}{-6}$ & 70.62 $\pm$ 0.09 & 63.27 $\pm$ 0.10 & $-7.34$, $1.9\mathrm{e}{-6}$ \\
California (CA) & 62.39 $\pm$ 0.13 & 64.07 $\pm$ 0.14 & $+1.68$, $2\mathrm{e}{-6}$ & 56.85 $\pm$ 0.09 & 47.35 $\pm$ 0.11 & $-9.50$, $1.9\mathrm{e}{-6}$ \\
Texas (TX)      & 63.11 $\pm$ 0.13 & 65.00 $\pm$ 0.11 & $+1.89$, $2\mathrm{e}{-6}$ & 59.44 $\pm$ 0.09 & 42.84 $\pm$ 0.14 & $-16.59$, $1.9\mathrm{e}{-6}$ \\
\midrule
Alabama (AL)    & 41.35 $\pm$ 0.17 & 40.57 $\pm$ 0.24 & $-0.78$, $0.036$ & 61.21 $\pm$ 0.18 & 50.17 $\pm$ 0.24 & $-11.04$, $1.9\mathrm{e}{-6}$ \\
Arizona (AZ)    & 43.90 $\pm$ 0.17 & 45.10 $\pm$ 0.19 & $+1.20$, $<\!1\mathrm{e}{-4}$ & 53.06 $\pm$ 0.15 & 40.78 $\pm$ 0.07 & $-12.27$, $1.9\mathrm{e}{-6}$ \\
\bottomrule
\end{tabular}
}
\end{table}


Next-region follows the cost-versus-gain pattern reported in the dense-prediction MTL line~\cite{vandenhende2022mtl}. MTL pays $-7.34$ pp (FL) to $-16.59$ pp (TX) on Acc@10 at every state, all statistically significant ($p\leq 1.9\mathrm{e}{-6}$). The pattern is broadly downward on AL\,$\to$\,AZ\,$\to$\,FL, with CA preserving the FL regime and TX breaking the monotonic shrink, plausibly under state-specific factors beyond raw class count (geographic POI distribution, user-density skew, per-state log-transition prior coverage); we do not read the AL\,$\to$\,FL contraction as an inferential scaling claim.

On Acc@10, the multi-task backbone trades the matched-head STAN-based ceiling for a regime that sits below a one-step Markov-1-region floor at four of five states (AL is the only state above the floor; per-state pp gaps: AL $+3.16$, AZ $-2.18$, FL $-1.78$, CA $-4.74$, TX $-12.10$). The matched-head Check2HGI ceiling sits $+5$ to $+14$ pp above Markov-1 at every state, so the cross-attention region head trades that ceiling for the joint regime, underperforming on transition structure a one-step Markov chain already captures. The $\Delta_m$ MRR axis at FL ($+2.33\,\%$, $n=25$, statistically significant in \S\ref{subsec:results-deltam}) characterises the joint tradeoff more cleanly than Acc@10 in this regime; we report both.

With the substrate fixed at the per-visit granularity, the multi-task backbone delivers a small category lift at four of five states and pays a sign-consistent cost on the harder head at every state, the classic multi-task tradeoff. The cost is structural to the cross-attention shared-backbone interaction with the region head; Section~\ref{sec:mechanism} reports drop-in replacements (FAMO, Aligned-MTL, hierarchical-additive-softmax region head) that fail to close the next-region gap.

\subsection{Joint $\Delta_m$: broadly negative; FL on MRR the exception}
\label{subsec:results-deltam}

On the joint $\Delta_m$ score~\cite{maninis2019attentive,vandenhende2022mtl} (the multi-task model's percent deviation from the matched-head ceilings averaged across tasks), FL is the only Pareto-positive cell on the primary axis at multi-seed; the other four states are negative on both axes. On the primary pair (category macro-F1, region MRR), FL carries an $n=25$ multi-seed sweep and reaches $\Delta_m\text{-MRR}=+2.33\,\%$ at $p=2.98\mathrm{e}{-8}$, 25/25 fold-pairs positive. The secondary pair at FL is $\Delta_m\text{-Acc@10}=-1.12\,\%$ at $p=3.20\mathrm{e}{-5}$ (4/25); MRR rewards near-correct candidates that raw top-10 accuracy misses, consistent with the Acc@10-saturated regime above. The other four states report at the $n=5$ Wilcoxon ceiling, with $\Delta_m$ negative on both axes (MRR from $-24.84\,\%$ to $-1.61\,\%$; Acc@10 from $-22.41\,\%$ to $-6.85\,\%$), at most $1/5$ fold-pairs positive.

\subsection{External baselines}
\label{subsec:results-external}

Against external baselines, the matched-head STL category ceiling on Check2HGI roughly doubles the strongest user-disjoint POI-RGNN reproduction, while the multi-task region row trails the strongest external region row. Table~\ref{tab:external} reports the comparison. On next-category, our matched-head STL ceiling on Check2HGI ($67.16$ at FL) sits well above our user-disjoint POI-RGNN reproduction~\cite{capanema2023poirgnn} ($34.49$); the published POI-RGNN evaluation used non-user-disjoint folds, so absolute numbers are not directly comparable. \S\ref{subsec:mechanism-pervisit} attributes ${\sim}72\%$ / ${\sim}64\%$ of the AL/AZ substrate gap specifically to per-visit context, so most of the lift is on the substrate axis, not from multi-task aggregation. On next-region, the matched-head STAN-based head with additive log-prior on the HGI substrate is the strongest external row at this scale; per-place HGI is a more natural fit for the region head, consistent with the substrate-axis reading. Our multi-task row trails this row by $-10$ to $-20$ pp on Acc@10, the substrate gap compounded by the cross-attention region cost.

\begin{table}[t]
\centering
\caption{External-baseline comparison at the headline scale (FL, CA, TX). The category block reports macro-F1; the region block reports Acc@10. Italicised STL rows are matched-head single-task ceilings on the Check2HGI substrate (\texttt{next\_gru} for category; \texttt{next\_stan\_flow} with the per-fold additive log-transition prior for region); the bold MTL row is our cross-attention multi-task model. POI-RGNN and MHA+PE numbers are our user-disjoint reproductions; the original POI-RGNN evaluation~\cite{capanema2022poirgnn} used non-user-disjoint folds, so the absolute values are not directly comparable. ReHDM and faithful STAN faithful runs are reported at AL/AZ/FL only; CA/TX cells are deferred to the camera-ready (the dual-level hypergraph in ReHDM scales quadratically with region cardinality, exceeding our compute budget at 8.5K and 6.5K regions respectively). Per-state cells are mean $\pm \sigma$ over the five user-disjoint folds.}
\label{tab:external}
\footnotesize
\resizebox{\textwidth}{!}{
\begin{tabular}{l r r r r r r}
\toprule
& \multicolumn{3}{c}{Next-category macro-F1} & \multicolumn{3}{c}{Next-region Acc@10} \\
\cmidrule(lr){2-4} \cmidrule(lr){5-7}
Method & FL & CA & TX & FL & CA & TX \\
\midrule
\multicolumn{7}{l}{\emph{External baselines (faithful reproductions)}} \\
\midrule
Markov-1 (floor)             & 23.98 $\pm$ 1.22 & 19.97 $\pm$ 1.77 & 18.17 $\pm$ 1.38 & 65.05 $\pm$ 0.93 & 52.09 $\pm$ 0.80 & 54.94 $\pm$ 0.46 \\
POI-RGNN \cite{capanema2022poirgnn} & 34.49 $\pm$ 1.14 & 31.78 $\pm$ 0.82 & 33.03 $\pm$ 0.70 & --              & --              & --              \\
MHA+PE \cite{zeng2019mhape}  & 32.06 $\pm$ 0.23 & 29.13 $\pm$ 0.71 & 29.91 $\pm$ 0.43 & --              & --              & --              \\
STAN \cite{luo2021stan} (HGI) & --              & --              & --              & 73.58 $\pm$ 0.43 & 60.45 $\pm$ 0.97 & 62.70 $\pm$ 0.37 \\
ReHDM \cite{li2025rehdm}     & --              & --              & --              & 65.68 $\pm$ 0.26 & --$^{\dagger}$  & --$^{\dagger}$  \\
\midrule
\multicolumn{7}{l}{\emph{Our matched-head single-task ceilings (Check2HGI substrate)}} \\
\midrule
\emph{STL \texttt{next\_gru} (cat ceiling)}        & 67.16 $\pm$ 0.13 & 62.39 $\pm$ 0.13 & 63.11 $\pm$ 0.13 & --              & --              & --              \\
\emph{STL \texttt{next\_stan\_flow} (reg ceiling)} & --              & --              & --              & 70.62 $\pm$ 0.09 & 56.85 $\pm$ 0.09 & 59.44 $\pm$ 0.09 \\
\midrule
\multicolumn{7}{l}{\emph{Proposed cross-attention multi-task model}} \\
\midrule
\textbf{MTL (ours)} & \textbf{68.56 $\pm$ 0.79} & \textbf{64.07 $\pm$ 0.14} & \textbf{65.00 $\pm$ 0.11} & 63.27 $\pm$ 0.10 & 47.35 $\pm$ 0.11 & 42.84 $\pm$ 0.14 \\
\bottomrule
\end{tabular}
}

\medskip
\footnotesize $^{\dagger}$ ReHDM at CA/TX deferred (camera-ready, compute budget).

\end{table}

